\documentclass[10pt,a4paper]{amsart}
\usepackage[margin=19mm]{geometry}
\usepackage{amsmath,amssymb,mathtools}
\usepackage[scr=rsfs,cal=euler]{mathalfa}
\usepackage{xfrac}
\usepackage[unicode,hidelinks,bookmarks=false]{hyperref}
\newcommand{\ie}{i.e.\ }
\newcommand{\cf}{cf.\ }
\newcommand{\resp}{respectively\ }
\newcommand{\Subsection}{\subsection}
\providecommand{\nobreakdash}{\nobreak\hbox{-}}
\setlength{\emergencystretch}{2em}
\multlinegap0pt
\usepackage{amssymb}
\usepackage{xcolor}
\newtheorem{theorem}{Theorem}
\title{The Jones polynomial of collections of open curves in $3-$space}
\author{Kasturi Barkataki, Eleni Panagiotou}
\date{Source: arXiv:2205.03474v1 (CC BY 4.0)}
\begin{document}
\maketitle

\noindent\emph{Source and licence.} The Jones polynomial of collections of open curves in $3-$space. Version arXiv:2205.03474v1 (CC BY 4.0), distributed by arXiv under the Creative Commons Attribution 4.0 International licence, http://creativecommons.org/licenses/by/4.0/, as recorded in the arXiv licence metadata for this exact version and shown on the abstract page https://arxiv.org/abs/2205.03474v1. Full licence notice: \texttt{licenses/arxiv-2205.03474-CC-BY-4.0.txt}.\par\smallskip

\noindent\textit{Editorial scope.} Counted unit: the article's theorem thm\_linktype\_jp (source0.tex lines 494--497). The article's complete original proof of it is delivered with it (source0.tex lines 499--529, one \texttt{proof} environment of 31 lines). No mathematics is written, repaired or re-derived: the statement and the proof are the article's own lines, in the article's own notation, and no retained line is substituted or re-flowed.  No further source-local context is needed: every cross-reference of every retained line resolves inside the two delivered spans.  Nothing was omitted from any retained span, so the package carries no omission record.  No retained line contains a citation, so the wrapper carries no bibliography and no bibliography entry.  No retained line contains a figure environment, an image-inclusion command or drawing code, so no image is delivered and the package declares no asset dependency.  Editorial formatting only, in addition: an AMS \texttt{amsart} preamble that restates the article's own declarations and macros used by the retained lines, namely the environments \texttt{theorem} on separate counters, together with the standard packages those lines need.  The retained environments print in this document as Theorem 1 in order of appearance, whereas the article prints the same items with the numbering of its own full document.  The complete native source of the article as distributed in the arXiv e-print for this exact version is bundled unchanged as \texttt{sources/125412/source0.tex}.
\begin{theorem}
\label{thm_linktype_jp}
Let $L$ be a link-type linkoid with $n$ components and $L_c$ be the corresponding link (the link that results from connecting the head to the leg of each component in a way that no new crossing is created). Then the Jones polynomials of $L$ and $L_c$ are equal, that is $f_L = f_{L_c}$.
\end{theorem}

\begin{proof}
By the definition of $L$ and $L_c$, we notice that $L$ can be created from $L_c$ by omitting $n$ arcs (closure arcs). These arcs connect the endpoints $2j-1$ and $2j$ for each component $l_{(2j-1,2j)}$ of $L$, where $j \in G=\{1, 2, \cdots, n\}$. Let us construct a decoration on $L_c$ by $L$, by keeping track of the endpoints of $L$ with labels. Thus $L_c$ is a union of $L$ with closure arcs. 

%To compare the states of $L_c$ to those of $L$, we will use a decoration on $L_c$, by labeling the endpoints, $G$, of the closure arcs on $L_c$. 

The closure arcs in $L$ do not create any new crossings, hence the total number of double points in $L_c$ is the same as that in $L$. Therefore, $Wr(L) = Wr(L_c)$. In the following, we show that the bracket polynomials, $\langle L \rangle$ and $\langle L_c \rangle$ are also equal. 


We know that the states of a link or a linkoid diagram are completely determined by the choice of smoothing at the crossings of the diagram.  Since $L$ and $L_c$ have identical crossings, for each state $S$ of the linkoid diagram $L$, there exists a state $S_c$ in the link diagram $L_c$, such that their smoothing labels are equal, that is, $\quad \sigma(S) = \sigma(S_c)$. 
By definition, the contribution of $S$ in the state sum expression of $\quad \langle L \rangle \quad$ is $ \quad A^{\sigma(S)} d^{|S|_{circ}  + |S|_{cyc}-1} \quad$, where $|S|_{circ}$ is the number of disjoint circles and 
$|S|_{cyc}$ is the number of segment cycles formed by the disjoint long segments in the state $S$.  
Similarly, the contribution of $S_c$ in the state sum expression of $\quad \langle L_c \rangle \quad$ is $\quad A^{\sigma(S)} d^{|S_c|_{circ} -1}. \quad$ We will prove that these two terms are equal by showing that $\quad |S|_{circ}+|S|_{cyc}= |S_c|_{circ}$. 

%In fact, the state $S$ of $L$ can be superposed on top of the state $S_c$ of $L_c$ such that every region of $S_c$ will be covered by $S$ except the $n$ closure arcs. Let us mark the endpoints of $L$ as decorations on $S_c$. Therefore, in the smoothed state, $S_c$, there will be some disjoint circles with decorations and some without decorations. 
%, every closure arc and the end points that it connects are intact in all states of $L_c$. Thus

Since the closure arcs of $L$ in $L_c$ do not introduce crossings, the state $S$ can be obtained from  $S_c$, by omitting the closure arcs. Thus, for a decorated $L_c$, $S_c$ will also be decorated with an even number of labels on each circle (since any closure arc has two endpoints and closure arcs are disjoint). Notice that if a circle in $S_c$ is not decorated, then it corresponds to a circle in $S$.
Let us express the total number of circles in $S_c$ as  $|S_c|_{circ} = |S_c|_{(circ,u)} + |S_c|_{(circ,d)} \quad ,$ where $\quad |S_c|_{(circ,u)}$ and $|S_c|_{(circ,d)} \quad$ denote the number of undecorated and decorated circles, respectively. Thus, $\quad |S|_{circ} = |S_c|_{(circ,u)}$. Therefore, we need to prove that $|S|_{cyc}= |S_c|_{(circ,d)}$.

%Note that if a decorated circle in a state $S_c$ contains an endpoint $x$ of a component of $L$, then it should also contain the other endpoint, say $x^\star$, of the same component because there exists an arc connecting the two endpoints which is unaffected by the smoothing. Therefore, any decorated circle should have even number of labels.

%Notice that the undecorated circles are the ones which do not involve any endpoints, therefore by superposition, for every undecorated circle in $S_c$ there exists a circle in $S$. Similarly, for every circle in $S$, there exists an undecorated circle in $S_c$. This implies, $\quad |S|_{circ} = |S_c|_{(circ,u)}$.

Once all the circles that do not involve endpoints are taken care of in both $S$ and $S_c$, we are left with $n$ long segments in $S$ and $|S_c|_{(circ,d)}$ number of decorated circles in $S_c$. 
Since $S$ is formed by $S_c$ by removing the closure arcs, a decorated circle in $S_c$ with $2k$ labels gives rise to $k$ long segments in $S$. 
We will show that these $k$ long segments form a segment cycle. 

Notice that for a decorated circle in $S_c$, two adjacent labels are either endpoints of a closure arc, or a new connection formed by the smoothings. So, each consecutive pair of labels in a decorated circle in $S_c$ that does not belong to a closure arc, it defines a pairing $J_S$ in $S$ and each consecutive pair of labels in the decorated circle in $S_c$ that belong to the same closure arc, they are related by $HL$. Then the corresponding $k$ long segments in $S$ define a segment cycle. %If a decorated circle in $S_c$ has only two labels, these correspond to the same segment, $i$, and after removing the closure arc, we obtain the segment $i$ of $S$, whose segment cycle has length 2. If a circle in $S_c$ contains more than 2 labels, by removing the closure arcs, we obtain segments in $S$ that connect endpoints of different initial segments. In general, after removing the closure arcs, we obtain the segments the connections between any two adjacent labels in this circle are alternately labelled as $J_S$ or $HL$. 

These $k$ long segments cannot overlap with any other decorated circle of $S_c$ because it will violate the fact that all the decorated circles of $S_c$ are disjoint. Therefore, for every decorated circle in $S_c$ there is a unique collection of long segments in $S$ such that these long segments form a segment cycle. This implies, $|S_c|_{(circ,d)} = |S|_{cyc}$. Therefore, $\langle L \rangle = \langle L_c \rangle$.  
\end{proof}

\end{document}
